\documentclass[11pt]{article}

\usepackage[margin=1in]{geometry}
\usepackage{amsmath,amssymb,amsthm}
\usepackage{bm}
\usepackage{hyperref}
\usepackage{parskip}

\renewcommand{\ne}{n_{e}}  % shadow LaTeX kernel's \ne ("not equal") -- not used in this doc
\newcommand{\nfc}{n_{\mathrm{FC}}}
\newcommand{\ncx}{n_{\mathrm{CX}}}
\newcommand{\Vfc}{\lvert V_{\mathrm{FC}}\rvert}
\newcommand{\Vcx}{\lvert V_{\mathrm{CX}}\rvert}
\newcommand{\Vfcr}{V_{\mathrm{FC},r}}
\newcommand{\Vcxr}{V_{\mathrm{CX},r}}
\newcommand{\Si}{S_{i}}
\newcommand{\Scx}{S_{\mathrm{CX}}}
\newcommand{\Sreact}{S}
\newcommand{\Dped}{D_{\mathrm{ped}}}
\newcommand{\fFC}{f_{\mathrm{FC}}}
\newcommand{\fCX}{f_{\mathrm{CX}}}
\newcommand{\fsa}[1]{\left\langle #1 \right\rangle}
\newcommand{\gradrsq}{\fsa{|\nabla r|^{2}}}
\newcommand{\neginf}{x=-\infty}
\newcommand{\dneneg}{\left.\dfrac{d\ne}{dx}\right|_{\neginf}}
\newcommand{\xin}{x_{\mathrm{inner}}}
\newcommand{\gFC}{\fsa{|\nabla r|^{2}\,\nfc}}
\newcommand{\gCX}{\fsa{|\nabla r|^{2}\,\ncx}}

\title{Derivation of Equation (16) in Saarelma \emph{et al.}~2023\\
  \large Nucl.\ Fusion \textbf{63} 052002 --- following the paper's procedure}
\author{}
\date{}

\begin{document}
\maketitle

This document derives Equation (16) of the Saarelma--Connor pedestal
neutral-penetration paper, starting from the three-fluid model Eqs.~(1)--(3),
flux-surface-averaging to get Eqs.~(4)--(6), introducing the form factors
Eq.~(7), and following the procedure of Eqs.~(8)--(10) through to~(16).
Every approximation is stated explicitly.

\section{Coordinates and flux-surface average}
\label{sec:coords}

We use straight-field-line coordinates $(r,\theta,\varphi)$ with Jacobian
$J = rR^{2}/R_{0}$ (Saarelma Sec.~2.1).  The flux-surface label
$r$ is the minor radius; $\theta$ is poloidal and $\varphi$ is toroidal.
In the narrow pedestal we introduce
\[
  x \;=\; r - r_{\mathrm{sep}}, \qquad r\simeq r_{\mathrm{sep}},
\]
so $x<0$ inside the pedestal and $x=0$ at the separatrix.

\paragraph{Flux-surface average.}
Following Eq.~(7) of the paper, the poloidal average is defined with
weight $R^{2}$:
\begin{equation}
  \fsa{A}(x) \;\equiv\; \frac{\oint R^{2}\,A\,d\theta}{\oint R^{2}\,d\theta}.
\end{equation}

\section{The three-fluid model (Eqs.~1--3)}
\label{sec:threefluid}

The paper models the edge plasma with three radially coupled fluids: the
electrons, the Franck--Condon (FC) neutrals produced by recycling and gas
puffing, and the warm charge-exchange (CX) neutrals produced when FC atoms
undergo a CX reaction with bulk ions:
\begin{align}
  \nabla\cdot(\Dped\,\nabla\ne)
    &= -\,\ne\,(\nfc + \ncx)\,\Si, \label{eq:1} \\
  \nabla\cdot(V_{FC}\,\nfc)
    &= -\,\ne\,\bigl(\nfc\Si + \nfc\Scx\bigr), \label{eq:2} \\
  \nabla\cdot(V_{CX}\,\ncx)
    &= -\,\ne\,\bigl(\ncx\Si - \tfrac{1}{2}\nfc\Scx\bigr). \label{eq:3}
\end{align}
The physics:
\begin{itemize}
  \item Eq.~\eqref{eq:1}: electron particle balance with diffusive flux
        $\Gamma_{e}=-\Dped\nabla\ne$ and an ionization source (only $\Si$
        appears; charge-exchange does not produce electrons).
  \item Eq.~\eqref{eq:2}: ballistic flow of FC neutrals at speed
        $\Vfc = \sqrt{8E_{FC}/\pi^{2}M_i}$ ($E_{FC}\simeq 3$\,eV) with loss
        to both ionization and CX.
  \item Eq.~\eqref{eq:3}: ballistic flow of the warm CX neutrals at
        $\Vcx = \sqrt{2T_i/\pi M_i}$, with loss to ionization and gain via
        the CX reaction on FC atoms.  The factor $\tfrac{1}{2}$ in front of
        $\nfc\Scx$ represents the half of the outgoing CX flux that is
        isotropised inward.
\end{itemize}
\textbf{Assumption A1 (locally mono-energetic neutrals):} Both neutral
populations are mono-energetic in velocity space (a $\delta$-function in
$\bm{v}$ at each spatial point), with a \emph{single inward radial
component}.  However, the two species behave differently in $x$:
\begin{itemize}
  \item The FC speed is set by Franck--Condon dissociation of $D_{2}$, a
        fixed atomic-physics quantity:
        \(
          \Vfc = \sqrt{8 E_{FC}/(\pi^{2} M_{i})}
        \)
        with $E_{FC}\approx 3$\,eV.  $\Vfc$ is genuinely constant in $x$
        (and in $\theta$).
  \item The CX neutral speed is set by the local ion thermal speed,
        \(
          \Vcx(x) = \sqrt{2 T_i(x)/(\pi M_{i})}.
        \)
        Since $T_i$ is a flux function (depends on $x$ but not $\theta$),
        $\Vcx$ varies \emph{across} the pedestal but not around a flux
        surface.
\end{itemize}
We therefore write $\Vfcr = -\Vfc$ (constant) and $\Vcxr = -\Vcx(x)$
(function of $x$).  The sign reflects that neutrals flow radially inward
from the separatrix.  Quasineutrality with $Z_{i}=1$ gives
$n_{i}\approx \ne$, which is why $\ne$ multiplies $\Scx$.

\section{Flux-surface integration (Eqs.~4--6)}
\label{sec:fsa-equations}

Applying $\oint R^{2}\,d\theta$ to Eq.~\eqref{eq:1}--\eqref{eq:3} and using
the identity
$\nabla\cdot\bm{F} = (1/J)\,\partial_{r}(J\,F^{r}) + \ldots$ in these
coordinates, the radial divergences integrate to
\begin{align}
  \frac{d}{dx}\!\left[\Dped\oint R^{2}|\nabla r|^{2}\,d\theta\,\frac{d\ne}{dx}\right]
    &= -\,\ne\oint R^{2}\,d\theta\,(\nfc+\ncx)\,\Si, \label{eq:4} \\
  \frac{d}{dx}\!\left[\oint R^{2}|\nabla r|^{2}\,d\theta\,\Vfcr\,\nfc\right]
    &= -\,\ne\oint R^{2}\,d\theta\,\bigl(\nfc\Si + \nfc\Scx\bigr), \label{eq:5} \\
  \frac{d}{dx}\!\left[\oint R^{2}|\nabla r|^{2}\,d\theta\,\Vcxr\,\ncx\right]
    &= -\,\ne\oint R^{2}\,d\theta\,\bigl(\ncx\Si - \tfrac{1}{2}\nfc\Scx\bigr). \label{eq:6}
\end{align}
Equations~\eqref{eq:4}--\eqref{eq:6} are Eqs.~(4)--(6) of the paper.  The
appearance of $|\nabla r|^{2}$ on the LHSs (but not on the RHSs) reflects
the fact that a radial flux couples to the surface area element via the
metric factor $|\nabla r|^{2}$, while the volumetric sources/sinks on the
RHSs are simply summed around the poloidal circumference weighted by
$R^{2}$.

\section{Form factors (Eq.~7)}
\label{sec:form-factors}

Define dimensionless form factors measuring the poloidal correlation
between $|\nabla r|^{2}$ and the neutral densities:
\begin{equation}
  \fFC(x) \;\equiv\; \frac{\fsa{|\nabla r|^{2}\,\nfc}}{\fsa{\nfc}\,\gradrsq},
  \qquad
  \fCX(x) \;\equiv\; \frac{\fsa{|\nabla r|^{2}\,\ncx}}{\fsa{\ncx}\,\gradrsq}.
  \label{eq:7}
\end{equation}
By construction $\fFC\to 1$ and $\fCX\to 1$ when the neutral density is
poloidally uniform on a flux surface.  In the opposite limit of a
neutral source localised at $\theta=\theta_{0}$,
$\fFC\simeq R^{2}(\theta_{0})|\nabla r|^{2}(\theta_{0})/\gradrsq$, so
$\fFC$ can be substantially different from unity.
The current Python implementation sets $\fFC=\fCX=1$
(\texttt{solver.py} line~520) --- i.e.\ neutrals assumed poloidally
symmetric.

\paragraph{Usage of the form factors.}
Rewriting~\eqref{eq:7}:
\begin{equation}
  \fsa{|\nabla r|^{2}\,\nfc} \;=\; \fFC\,\fsa{\nfc}\,\gradrsq,
  \qquad
  \fsa{|\nabla r|^{2}\,\ncx} \;=\; \fCX\,\fsa{\ncx}\,\gradrsq.
  \label{eq:7b}
\end{equation}

\section{Reduction of Eq.~(5) to Eq.~(8)}
\label{sec:eq-8}

Divide~\eqref{eq:5} by the slowly-varying geometric factor
$\oint R^{2}\,d\theta$ and pull out the constant $\Vfcr=-\Vfc$:
\begin{equation}
  -\,\Vfc\,\frac{d}{dx}\fsa{|\nabla r|^{2}\,\nfc}
    \;=\; -\,\ne\,(\Si+\Scx)\,\fsa{\nfc}.
\end{equation}
Insert~\eqref{eq:7b} for $\fsa{|\nabla r|^{2}\nfc}$:
\begin{equation}
  \Vfc\,\frac{d}{dx}\!\bigl[\fFC\,\fsa{\nfc}\,\gradrsq\bigr]
    \;=\; \ne\,(\Si+\Scx)\,\fsa{\nfc}.
\end{equation}

\textbf{Assumption A2 (slowly varying geometry and form factor):} the
geometric flux-surface integral $\gradrsq$ and the form factor $\fFC$ vary
slowly with $x$ compared to $\fsa{\nfc}$ across the narrow pedestal, and
are treated as constants in $d/dx$.  This gives
\begin{equation}
  \Vfc\,\fFC\,\gradrsq\,\frac{d\fsa{\nfc}}{dx}
    \;=\; \ne\,(\Si+\Scx)\,\fsa{\nfc}.
  \label{eq:8-ours}
\end{equation}

\textbf{Assumption A3 (slab-like pedestal for the neutrals):} the paper now
drops the explicit factor $\gradrsq$ on the LHS of~\eqref{eq:8-ours}.  This
is equivalent to absorbing $\gradrsq$ into the neutral flow (i.e.\ treating
the pedestal as slab-like for the \emph{neutral} kinetic balance, while
retaining it in the electron transport equation~\eqref{eq:4}).  With
$1+\Scx/\Si = (\Si+\Scx)/\Si$ this leaves
\begin{equation}
  \boxed{\;
    \ne\Si\fsa{\nfc} \;=\; \frac{\Vfc\,\fFC}{1 + \Scx/\Si}\,\frac{d\fsa{\nfc}}{dx}.
  \;}
  \label{eq:8}
\end{equation}
Equation~\eqref{eq:8} is Eq.~(8) of the paper.  For the code's
$\fFC\to 1$ limit it reduces to the simple FC-kinetic balance
$\Vfc\,d\fsa{\nfc}/dx = \ne(\Si+\Scx)\fsa{\nfc}$.

\section{Reduction of Eq.~(6) to Eq.~(9)}
\label{sec:eq-9}

Apply the same steps to~\eqref{eq:6}.  One subtlety: unlike $\Vfcr$, the
radial component $\Vcxr(x) = -\Vcx(x)$ depends on $x$ through
$T_i(x)$ (see Assumption~A1).  So pulling $\Vcxr$ outside of $d/dx$
drops the term $(d\Vcx/dx)\,\fsa{|\nabla r|^{2}\ncx}\,\oint R^{2}d\theta$:
\begin{align}
  \frac{d}{dx}\!\left[\oint R^{2}\,d\theta\,\Vcxr(x)\,\fsa{|\nabla r|^{2}\ncx}\right]
  &= \Vcxr(x)\,\frac{d}{dx}\!\left[\oint R^{2}d\theta\,\fsa{|\nabla r|^{2}\ncx}\right] \notag\\
  &\qquad + \frac{d\Vcxr}{dx}\,\oint R^{2}d\theta\,\fsa{|\nabla r|^{2}\ncx}.
\end{align}
\textbf{Assumption A2b (slowly varying CX speed):} drop the
$d\Vcx/dx$ term, i.e.\ treat $\Vcx$ as constant over the scale on which
$\ncx$ varies.  This is a standard ``$T_i$ varies slowly over a neutral
mean-free path'' approximation.  Quantitatively,
$\tfrac{1}{\Vcx}\tfrac{d\Vcx}{dx}=\tfrac{1}{2T_i}\tfrac{dT_i}{dx}$, so the
error is $(L_{T_i}/L_{n_{\mathrm{CX}}})^{-1}$, which is small provided
$\ncx$ decays faster than $T_i$ drops --- typical in the pedestal.

With this approximation and after using~\eqref{eq:8} for the CX-source
term on the RHS, we obtain
\begin{equation}
  \ne\Si\fsa{\ncx} \;=\; \Vcx(x)\,\fCX\,\frac{d\fsa{\ncx}}{dx}
    \;+\; \frac{\Vfc\,\fFC\,\Scx}{2(\Si+\Scx)}\,\frac{d\fsa{\nfc}}{dx}.
  \label{eq:9}
\end{equation}
This is Eq.~(9) of the paper.  The same geometry/form-factor slow-variation
and slab-like assumptions (A2, A3) have been made, plus the new A2b on
$\Vcx$.  When $\Vcx(x)$ appears in the code it is evaluated
\emph{locally} at each flux surface from $T_i(x)$ (see
\texttt{solver.py} definitions of \texttt{V\_CX}), consistent with A2b:
$\Vcx$ is mono-energetic on each flux surface but its value follows
$T_i(x)$ --- it is only the $d\Vcx/dx$ piece of the divergence that is
dropped.

\section{First integration: Eq.~(10)}
\label{sec:eq-10}

Divide~\eqref{eq:4} by $\oint R^{2}\,d\theta$:
\begin{equation}
  \frac{d}{dx}\!\left[\gradrsq\,\Dped\,\frac{d\ne}{dx}\right]
    \;=\; -\,\ne\,\Si\,\bigl(\fsa{\nfc}+\fsa{\ncx}\bigr).
  \label{eq:4-1D}
\end{equation}
Sum the RHSs of~\eqref{eq:8} and~\eqref{eq:9}:
\begin{align}
  \ne\Si\bigl(\fsa{\nfc}+\fsa{\ncx}\bigr)
    &= \underbrace{\left[\frac{1}{1+\Scx/\Si} + \frac{\Scx}{2(\Si+\Scx)}\right]}_{\displaystyle\text{collect}}\Vfc\,\fFC\,\frac{d\fsa{\nfc}}{dx}
      + \Vcx\,\fCX\,\frac{d\fsa{\ncx}}{dx}.
\end{align}
The bracketed coefficient simplifies:
\begin{equation}
  \frac{\Si}{\Si+\Scx} + \frac{\Scx}{2(\Si+\Scx)}
  \;=\; \frac{2\Si+\Scx}{2(\Si+\Scx)}
  \;=\; \frac{\Si+\Scx/2}{\Si+\Scx}
  \;=\; \frac{1+\Scx/(2\Si)}{1+\Scx/\Si}.
\end{equation}
Substitute into~\eqref{eq:4-1D} and integrate once in $x$.  Using A2 for
$\Vfc\fFC$ (exact, since $\Vfc$ is constant and $\fFC$ is slowly
varying) and A2b together with A2 for $\Vcx\fCX$ (so both pass through
the $d/dx$), we obtain
\begin{equation}
  \boxed{\;
    \gradrsq\,\Dped\,\frac{d\ne}{dx}
    \;=\;
    -\,\frac{1+\Scx/(2\Si)}{1+\Scx/\Si}\,\Vfc\,\fFC\,\fsa{\nfc}
    \;-\; \Vcx\,\fCX\,\fsa{\ncx}
    \;+\; C.
  \;}
  \label{eq:10}
\end{equation}
The integration constant $C$ is fixed by requiring the pedestal to match
the core: as $x\to-\infty$, both $\fsa{\nfc}\to 0$ and $\fsa{\ncx}\to 0$,
so
\begin{equation}
  C \;=\; \left[\gradrsq\,\Dped\,\frac{d\ne}{dx}\right]_{\neginf}
  \;\simeq\; \gradrsq\,\Dped\,\dneneg.
  \label{eq:Cvalue}
\end{equation}
In Saarelma~[5] $C$ is taken to zero; here we keep it finite to allow a
non-zero density gradient at the inner boundary (as the code does via the
BC $n_{e}'(\xin)=\dneneg$).

\section{The no-CX limit: derivation of Eq.~(16)}
\label{sec:eq-16}

To obtain the starting point for the iterative solve of the full
equation~(15), Saarelma takes $\fsa{\ncx}\to 0$ in Eq.~\eqref{eq:10}:
\begin{equation}
  \gradrsq\,\Dped\,\frac{d\ne}{dx}
    \;=\; -\,\frac{1+\Scx/(2\Si)}{1+\Scx/\Si}\,\Vfc\,\fFC\,\fsa{\nfc} + C.
  \label{eq:10-noCX}
\end{equation}
Solve for $\fsa{\nfc}$:
\begin{equation}
  \fsa{\nfc}
    \;=\; \frac{1+\Scx/\Si}{1+\Scx/(2\Si)}\;
          \frac{C - \gradrsq\,\Dped\,\dfrac{d\ne}{dx}}{\Vfc\,\fFC}.
  \label{eq:nfc-closure}
\end{equation}
Substitute~\eqref{eq:Cvalue} for $C$ and substitute~\eqref{eq:nfc-closure}
back into~\eqref{eq:4-1D} (which, in the no-CX limit, reads
$\tfrac{d}{dx}[\gradrsq\Dped\ne'] = -\ne\Si\fsa{\nfc}$):
\begin{align}
  \frac{d}{dx}\!\left[\gradrsq\,\Dped\,\frac{d\ne}{dx}\right]
    &= -\,\ne\,\Si\,\frac{1+\Scx/\Si}{1+\Scx/(2\Si)}\,
         \frac{\gradrsq\,\Dped\,\bigl(\dneneg - d\ne/dx\bigr)}{\Vfc\,\fFC} \\[4pt]
    &= \frac{\ne\,(\Si+\Scx)}{1+\Scx/(2\Si)}\,
       \frac{\gradrsq\,\Dped}{\Vfc\,\fFC}\,\bigl(\tfrac{d\ne}{dx} - \dneneg\bigr).
  \label{eq:16-full}
\end{align}

Equation~\eqref{eq:16-full} is our \emph{complete} form of Eq.~(16),
keeping every geometric factor and every form factor explicit.

\paragraph{Saarelma's further simplifications to obtain Eq.~(16) as written
in the paper.}
\begin{itemize}
  \item \textbf{Assumption A4 ($\fFC=1$):} poloidal symmetry of FC
        neutrals.  The code sets \texttt{self.fFC = 1} via
        \texttt{form\_factor()}.
  \item \textbf{Assumption A5 (drop the shape factor):} approximate
        $(1+\Scx/(2\Si))/1 \approx 1$, i.e., $\Scx/(2\Si)\ll 1$.  This is
        the weakest assumption --- it is not strictly correct in the
        pedestal, but is taken in the paper to write~\eqref{eq:16-full} in
        its compact form.
  \item \textbf{Retain $\gradrsq$ only in the transport operator.}  The
        numerator $\gradrsq\,\Dped$ on the RHS of~\eqref{eq:16-full} is
        treated as the prefactor $\Dped$ times a factor of $\gradrsq$ that
        cancels symbolically with the $\gradrsq$ already on the LHS of the
        transport operator, leaving simply $\Dped$ on the RHS.  (This
        cancellation is symbolic rather than rigorous --- it amounts to an
        additional $\gradrsq\approx 1$ approximation inside the neutral
        closure.)
\end{itemize}

Combining A4 and A5 with the cancellation above, Eq.~\eqref{eq:16-full}
reduces to the form written in the paper:
\begin{equation}
  \boxed{\;
    \frac{d}{dx}\!\left[\gradrsq\,\Dped\,\frac{d\ne}{dx}\right]
    \;=\; \ne\,\Dped\,\frac{\Si+\Scx}{\Vfc}\,
          \left(\frac{d\ne}{dx} - \dneneg\right).
  \;}
  \label{eq:16}
\end{equation}
This is Eq.~(16) of Saarelma \emph{et al.}~2023.  Matches
\texttt{first\_step} in \texttt{solver.py} exactly.

\section{Physical form of the ODE}
\label{sec:phys-ode}

Define $f(x)\equiv \gradrsq\,\Dped$ (this is the \texttt{f\_arr} in the
code, \emph{not} the form factor $\fFC$).  Expand the LHS
of~\eqref{eq:16}:
\begin{equation}
  f(x)\,\ne'' + f'(x)\,\ne'
    \;=\; \ne\,\Dped\,\frac{\Si+\Scx}{\Vfc}\,(\ne' - \dneneg).
\end{equation}
Divide by $f$ and rearrange:
\begin{equation}
  \boxed{\;
    \ne'' \;=\; c_A(x)\,\ne\,\ne' \;-\; c_B(x)\,\ne \;-\; c_K(x)\,\ne',
  \;}
  \label{eq:phys-ODE}
\end{equation}
with physical coefficients
\begin{align}
  c_A(x) &\;=\; \frac{\Dped(\Si+\Scx)}{\Vfc\,f}
           \;=\; \frac{\Si+\Scx}{\Vfc\,\gradrsq}, \\
  c_B(x) &\;=\; \frac{\Dped(\Si+\Scx)\,\dneneg}{\Vfc\,f}
           \;=\; \frac{(\Si+\Scx)\,\dneneg}{\Vfc\,\gradrsq}, \\
  c_K(x) &\;=\; \frac{f'(x)}{f(x)}.
\end{align}
These are the coefficients of Eq.~(16) in physical space.

\section{Non-dimensionalization}
\label{sec:nondim}

Introduce a density scale $n_{0}$ and a length scale $L$:
\begin{equation}
  N(\xi) \;\equiv\; \frac{\ne(x)}{n_{0}}, \qquad \xi \;\equiv\; \frac{x}{L}.
\end{equation}
In the code $n_{0}=\texttt{ne\_x0}=\ne(0)$ and $L=|\xin|$
(see \texttt{non\_dimensionalize}).  The chain rule gives
\begin{equation}
  \ne' = \frac{n_{0}}{L}\,N', \qquad \ne'' = \frac{n_{0}}{L^{2}}\,N'',
\end{equation}
and the non-dimensional Neumann BC at $\xi=-1$ is
\begin{equation}
  \left.\frac{dN}{d\xi}\right|_{\xi=-1} \;=\; \frac{L}{n_{0}}\,\dneneg.
\end{equation}
Substitute term by term into~\eqref{eq:phys-ODE}:
\begin{align}
  \ne'' &= \frac{n_{0}}{L^{2}}\,N'', \\
  c_A\,\ne\,\ne' &= c_A\,(n_{0}N)\frac{n_{0}}{L}N' = \frac{n_{0}^{2}}{L}\,c_A\,N\,N', \\
  c_B\,\ne &= n_{0}\,c_B\,N, \\
  c_K\,\ne' &= \frac{n_{0}}{L}\,c_K\,N'.
\end{align}
Multiply the substituted equation by $L^{2}/n_{0}$:
\begin{equation}
  N'' \;=\; \underbrace{(n_{0}L\,c_A)}_{A}\,N\,N'
         \;-\; \underbrace{(L^{2}\,c_B)}_{B}\,N
         \;-\; \underbrace{(L\,c_K)}_{K}\,N'.
\end{equation}
Plugging in the physical coefficients gives
\begin{equation}
  \boxed{\;
    \begin{aligned}
      A \;&=\; \frac{n_{0}\,L\,\Dped\,(\Si+\Scx)}{\Vfc\,f}
           \;=\; \frac{n_{0}\,L\,(\Si+\Scx)}{\Vfc\,\gradrsq}, \\[4pt]
      B \;&=\; \frac{L^{2}\,\dneneg\,\Dped\,(\Si+\Scx)}{\Vfc\,f}
           \;=\; \frac{n_{0}\,L\,(L/n_{0})\dneneg\,\Dped\,(\Si+\Scx)}{\Vfc\,f}, \\[4pt]
      K \;&=\; \frac{L\,f'(x)}{f(x)} \;=\; \frac{1}{f}\,\frac{df}{d\xi}.
    \end{aligned}
  \;}
\end{equation}
The two forms of $B$ are equivalent because
$\dneneg = (n_{0}/L)\bigl(dN/d\xi\bigr)_{\xi=-1}$, so
$L^{2}\dneneg = n_{0}L\bigl(dN/d\xi\bigr)_{\xi=-1}$; the code writes $B$ in
this second form (see the ODE body in \texttt{first\_step}).  These are
precisely the coefficients implemented in lines~760--762 of
\texttt{solver.py}.

\section{Dimensional check}
\label{sec:dimcheck}

With $[\Dped]=[f]=\mathrm{m^{2}/s}$, $[\Sreact]=\mathrm{m^{3}/s}$,
$[\Vfc]=\mathrm{m/s}$, $[n_{0}]=\mathrm{m^{-3}}$, $[L]=\mathrm{m}$,
$[\dneneg]=\mathrm{m^{-4}}$:
\begin{align}
  [A] \;&=\; \frac{\mathrm{m^{-3}}\cdot\mathrm{m}\cdot\tfrac{\mathrm{m^{2}}}{\mathrm{s}}\cdot\tfrac{\mathrm{m^{3}}}{\mathrm{s}}}
                   {\tfrac{\mathrm{m}}{\mathrm{s}}\cdot\tfrac{\mathrm{m^{2}}}{\mathrm{s}}}
               \;=\; 1, \\[4pt]
  [B] \;&=\; \frac{\mathrm{m^{2}}\cdot\mathrm{m^{-4}}\cdot\tfrac{\mathrm{m^{2}}}{\mathrm{s}}\cdot\tfrac{\mathrm{m^{3}}}{\mathrm{s}}}
                   {\tfrac{\mathrm{m}}{\mathrm{s}}\cdot\tfrac{\mathrm{m^{2}}}{\mathrm{s}}}
               \;=\; 1, \\[4pt]
  [K] \;&=\; \mathrm{m}\cdot\frac{(\mathrm{m^{2}/s})/\mathrm{m}}{\mathrm{m^{2}/s}}
               \;=\; 1.
\end{align}
All three coefficients are dimensionless.  \hfill$\checkmark$

\section{Summary of assumptions}
\label{sec:assumptions}

For easy cross-reference when comparing to the code, the approximations
stacked in the derivation are:

\begin{itemize}
  \item[\textbf{A1}] \textbf{Locally mono-energetic neutrals.} Both
        species have a single radial speed at each point, with magnitudes
        $\Vfc=\sqrt{8E_{FC}/(\pi^{2}M_{i})}$ (genuinely constant, atomic
        physics) and $\Vcx(x)=\sqrt{2T_{i}(x)/(\pi M_{i})}$ (constant
        around a flux surface but varying across the pedestal through
        $T_i(x)$).
  \item[\textbf{A2}] \textbf{Slowly varying geometry.} $\gradrsq$ and
        $\oint R^{2}\,d\theta$ are approximately constant in $x$ across
        the pedestal width.  Same for the form factors $\fFC$, $\fCX$.
  \item[\textbf{A2b}] \textbf{Slowly varying $\Vcx$.} When going from
        Eq.~\eqref{eq:6}$\to$\eqref{eq:9}, the $(d\Vcx/dx)\,\fsa{\ncx}$
        term is dropped.  Required because $\Vcx\propto\sqrt{T_i(x)}$
        truly varies across the pedestal, unlike $\Vfc$.  Valid when
        $\ncx$ decays faster than $T_i$ drops, i.e.\
        $L_{n_{\mathrm{CX}}}\ll L_{T_i}$.
  \item[\textbf{A3}] \textbf{Slab-like neutral closure.} In going from
        Eq.~\eqref{eq:5}$\to$\eqref{eq:8} and \eqref{eq:6}$\to$\eqref{eq:9},
        an explicit $\gradrsq$ factor is dropped, equivalent to using a
        slab-like closure for the neutral transport equation.  (The
        transport operator on the LHS of~\eqref{eq:4} retains its
        $\gradrsq$.)
  \item[\textbf{A4}] \textbf{Poloidal symmetry of FC neutrals,
        $\fFC=1$.} Assumed when collapsing Eq.~\eqref{eq:16-full} to the
        paper's Eq.~\eqref{eq:16}.  Set in
        \texttt{solver.py} line~520 by \texttt{form\_factor()}.
  \item[\textbf{A5}] \textbf{Weak CX contribution in FC balance.}
        $1+\Scx/(2\Si)\approx 1$, dropped when going from
        Eq.~\eqref{eq:16-full} to~\eqref{eq:16}.
  \item[\textbf{A6}] \textbf{No CX neutrals.} $\fsa{\ncx}=0$.  This is the
        ``first-step'' limit used in \texttt{first\_step()}; lifted
        iteratively in \texttt{solve()}.
  \item[\textbf{A7}] \textbf{Boundary condition at $x_{\mathrm{inner}}$.}
        $\ne'(\xin)=\dneneg$ (Neumann, finite) and
        $\fsa{\nfc}(\xin)\approx 0$.  Allows a non-zero $C$ in
        Eq.~\eqref{eq:10}, unlike~[5] where $C\equiv 0$.
  \item[\textbf{A8}] \textbf{Quasineutrality with $Z_{i}=1$.}
        $n_{i}\approx\ne$, used implicitly when $\ne$ multiplies $\Scx$ in
        Eqs.~\eqref{eq:2}--\eqref{eq:3}.
\end{itemize}

\end{document}
